\section{Experimental setup}
\label{sec:experimental-setup}

We instantiate \method on three benchmarks that share the same
architecture, objective, and training loop (\Cref{sec:method}) and
differ only in the observation, the state and action spaces, and the
hand-derived solver $S$. This section gives the exact solvers
(\Cref{sec:setup-benchmarks}), the diagnostic that checks each solver
is expressive enough to explain real transitions before any claim
about the probe is made (\Cref{sec:setup-adequacy}), the data and
architecture (\Cref{sec:setup-data}), and the protocol used to produce
the results in the next section (\Cref{sec:setup-protocol}).

\subsection{Benchmarks and physics solvers}
\label{sec:setup-benchmarks}

\paragraph{\textsc{PokeWorld}.} A synthetic point-mass benchmark built
for this study. A poker, driven by the action, pushes a disc object
subject to contact and viscous drag. State is $(x, y, v_x, v_y)$,
action is the poker's $(x, y)$ position, and the physics is
\begin{equation}
  m \, \dot v = k \cdot \mathrm{overlap}(x, x_{\mathrm{poker}}) \cdot
    \hat n - c\, v,
\end{equation}
where $\mathrm{overlap}$ is the signed penetration between poker and
object radii (zero outside contact) and $\hat n$ the contact normal.
$\theta = (m, k, c, r_{\mathrm{poker}}, r_{\mathrm{object}})$, five
parameters, is sampled independently per episode and rendered as two
soft discs on a blank field, so mass, stiffness, and drag are hidden
behind identical-looking rollouts and the only route to them is the
dynamics. Because only the ratios $k/m$ and $c/m$ affect the
trajectory, $m$, $k$, and $c$ are not individually identifiable from
state transitions alone; this is a property of the dynamics, not a
limitation of the probe, and \Cref{sec:setup-protocol} reports it
rather than working around it. \textsc{PokeWorld} is the one benchmark
where $\theta$ is actually known, which is what makes the $R^2$
certificate possible.

\paragraph{Cart-pole.} State is \texttt{dm\_control}'s
$(x, \phi, \dot x, \dot\phi)$ (cart position, pole angle from upright,
their velocities), action is a scalar force, and the solver is the
standard cart-pole equations of motion with a uniform-rod pole (the
$4/3$ term is its moment of inertia about the hinge):
\begin{equation}
  \tau = \frac{F + m_p \ell \dot\phi^2 \sin\phi - b\, \dot x}{m_c +
    m_p}, \qquad
  \ddot\phi = \frac{g \sin\phi - \tau \cos\phi}
    {\ell\!\left(\tfrac{4}{3} - \tfrac{m_p \cos^2\phi}{m_c + m_p}\right)},
  \qquad
  \ddot x = \tau - \frac{m_p \ell \ddot\phi \cos\phi}{m_c + m_p},
\end{equation}
integrated with semi-implicit Euler, matching MuJoCo's own integration
style. $\theta = (m_c, m_p, \ell, g, k_F, b)$: cart mass, pole mass,
pole half-length, gravity, the force-to-actuation gain, and cart
damping. Six parameters, none of them labeled in the data. Episodes
are live rollouts of \texttt{swm/CartpoleDMControl-v0}.

\paragraph{Push-T.} State is the environment's own seven-dimensional
observation, agent position, block position, block angle, and the
agent's (not the block's) velocity, a naming quirk in the environment
worth flagging because it is easy to get backward when writing a
solver against it. The agent follows the environment's exact PD law,
$a = k_p(\mathrm{target} - x) - k_v \dot x$, so its two gains are
genuinely identifiable rather than merely fit. The block has no
directly observed velocity, so we model it quasi-statically: contact
is a soft penetration spring between the agent and a disc standing in
for the T-shaped block, and the resulting force and torque map to
block motion through an isotropic limit-surface mobility, with no
block momentum. Rotation exists only because of a body-fixed
center-of-friction offset from the block's geometric center; a disc
pushed through its own center has zero lever arm and would never spin,
so this offset is what makes the angular channel identifiable at all.
$\theta = (k_p, k_v, r_{\mathrm{agent}}, r_{\mathrm{block}}, k,
\mu_{\mathrm{lin}}, \mu_{\mathrm{ang}}, c_x, c_y)$, nine parameters.
Approximating a T-block as a disc is a known structural mismatch,
documented and measured rather than hidden (\Cref{sec:setup-adequacy}).
Episodes are live rollouts of \texttt{swm/PushT-v1} under a scripted
policy that steers the agent toward the block; uniform random actions
almost never make contact, which would leave the block's physics
completely unexercised in the data.

All three solvers are hand-derived, hardcoded integrators written as
plain tensor operations, differentiable only in the sense that any
such computation is differentiable under automatic differentiation,
not through a general-purpose differentiable physics engine. $\theta$
is mapped from an unconstrained probe output into each solver's
physically valid box via $\mathrm{lo} + (\mathrm{hi} - \mathrm{lo})
\cdot \sigma(\theta_{\mathrm{raw}})$, with $\mathrm{lo}$, $\mathrm{hi}$
fixed per parameter (e.g.\ cart mass in $[0.1, 5.0]\,\mathrm{kg}$,
gravity in $[1, 20]\,\mathrm{m/s^2}$); \Cref{sec:method-forward} gives
the general form.

\subsection{Solver adequacy}
\label{sec:setup-adequacy}

Before asking whether the probe can recover $\theta$ from the latent,
we first ask a prior question: is each solver's fixed functional form
even capable of explaining real transitions, for some $\theta$? A null
result on the probe would be uninformative without this check, since
it could mean either that the physics is not in the latent or that the
physics is not in the solver. We answer it by fitting a free
per-episode $\theta$ directly against real transitions with gradient
descent, an oracle that exists only for this diagnostic and is never
used during \method training, where $\theta$ is always a probe output
(\Cref{sec:method-invariants}). \Cref{tab:solver-adequacy} reports
scaled RMSE against a persistence baseline ($s_{t+1} = s_t$) and
against the solver run at its nominal, un-fit $\theta$.

\begin{table}[t]
\centering
\begin{tabular}{lccc}
\toprule
Benchmark & Persistence & Nominal $\theta$ & Fitted $\theta$ (oracle) \\
\midrule
\textsc{PokeWorld} & 0.312 & 0.228 & \textbf{0.0001} \\
Cart-pole          & 0.166 & 0.006 & \textbf{0.0025} \\
Push-T             & 0.352 & 0.031 & \textbf{0.021}  \\
\bottomrule
\end{tabular}
\caption{Solver adequacy, scaled RMSE, lower is better, produced by
\texttt{scripts/smoke/validate\_solvers.py}. All three solvers fit
real transitions to near zero error given an oracle per-episode
$\theta$, so a low $R^2$ from the trained probe would indicate a
latent that has not internalized the physics, not a solver that
cannot express it. On Push-T, the
disc approximation of the T-block is the dominant source of remaining
error even at the oracle fit; \Cref{sec:setup-benchmarks} documents
why.}
\label{tab:solver-adequacy}
\end{table}

Cart-pole's fitted $\theta$ recovers physically correct parameters
from the MuJoCo rollouts (gravity $\approx 9.81\,\mathrm{m/s^2}$, cart
mass $\approx 1.01\,\mathrm{kg}$, pole half-length $\approx
0.517\,\mathrm{m}$), and Push-T's fit recovers the environment's exact
controller gains ($k_p{=}100$, $k_v{=}20$), the sanity check that the
agent half of that solver is exactly correct, not merely adequate.

\subsection{Data and architecture}
\label{sec:setup-data}

\paragraph{Data.} \textsc{PokeWorld} episodes are generated on the
fly (64 training, 16 validation, disjoint random seeds, 64 steps each).
Cart-pole and Push-T episodes are live environment rollouts (64
training / 16 validation episodes of 64 steps for cart-pole; 32
training / 8 validation episodes of 48 steps for Push-T), rendered at
$64\times64$. Validation episodes are drawn from an entirely separate
seed range rather than held out from the same rollouts, so no window
in validation shares a trajectory with training. Episodes are sliced
into overlapping windows of $T{=}8$ frames with stride 1; a window
never straddles an episode boundary, since $\theta$ is constant within
one (\Cref{sec:method-forward}).

\paragraph{Encoder.} Two options, swappable from config: a small
dependency-free CNN producing a $4\times4$ grid of patch tokens
(the default here, since it needs no pretrained weights and keeps the
full study runnable on a single consumer GPU), or a frozen DINOv2-small
backbone for experiments that need the real pretrained representation
this paper's motivation is about. A state-only encoder is also
available as an ablation that removes pixels and the visual
representation entirely, isolating the two-path objective from
whatever the encoder contributes.

\paragraph{Predictor and heads.} \pathA's predictor is a 4-layer,
8-head block-causal transformer (hidden size 256, MLP width 512) over
patch tokens, following DINO-WM \citep{zhou2024dinowm}. The state
decoder is a 2-layer MLP over mean-pooled tokens. \pathB's probe is,
by default, a single linear map over mean-pooled, time-pooled tokens
(zero hidden layers), initialized with small weights so the initial
$\theta$ starts near each solver's nominal parameter set rather than
at a random point in the box.

\paragraph{Optimization.} AdamW, learning rate $3\times10^{-4}$,
weight decay $10^{-4}$, batch size 32, 30 epochs, linear warmup over
the first 5\% of steps followed by cosine decay, gradient norm clipped
to 1.0. State and action normalizers are z-scored from 512 sampled
training windows once, before training, and then frozen as buffers.
$\alpha{=}1$, $\beta{=}0$ in \Cref{eq:loss} unless stated otherwise.

\subsection{Evaluation protocol}
\label{sec:setup-protocol}

Three measurements test the central claim from different angles: that
$\theta$ is functionally responsible for what the model predicts, not
merely present in the latent somewhere.

\paragraph{Identifiability certificate.} On \textsc{PokeWorld}, where
true $\theta$ is known, we accumulate predicted and true $\theta$ over
a full validation epoch (never a single minibatch, since $R^2$ is only
meaningful once the true parameter actually varies across the sample)
and report per-parameter $R^2 = 1 - \mathrm{SS}_{\mathrm{res}} /
\mathrm{SS}_{\mathrm{tot}}$. Parameters that are constant by
construction, or constant within a given sample, return an explicit
undefined value rather than a numerically large but meaningless score,
and we report the ratios $k/m$ and $c/m$ alongside the raw parameters
given the non-identifiability in \Cref{sec:setup-benchmarks}.

\paragraph{Out-of-distribution transfer.} We evaluate a trained
checkpoint on rollouts whose true physical parameters lie outside the
training range (for example, an object mass or a controller gain
beyond the box the model was trained on) and compare \pathA's
free-running prediction against \pathB's. A latent that only memorized
in-distribution correlations, rather than a $\theta$ doing physical
work, should show \pathA drifting while \pathB, constrained by the
solver's fixed equations regardless of what $z$ says, continues to
track the true dynamics.

\paragraph{Downstream control.} We plug a trained \method checkpoint
into the repository's model-predictive control stack (CEM/MPPI over
the learned dynamics) and measure goal-reaching success rate under the
same out-of-distribution parameter shift, against a predictive-only
DINO-WM baseline trained identically but without \pathB. This is the
test that a merely decodable but functionally inert $\theta$ would
fail: decoding a physical parameter from a latent is not the same
claim as that parameter changing what the plan actually does.
